% =========================================================================
% PLANTILLA MAESTRA: APUNTES TÉCNICOS COMPACTOS (ESTILO PASTEL)
% =========================================================================
% Características:
% - Márgenes reducidos (1.25 cm) y tipografía compacta (9pt con lmodern).
% - Sin encabezado superior, sin portada separada y sin tabla de contenidos.
% - Banner superior estilizado y compacto en la página inicial.
% - Paleta temática de 4 colores pastel para bloques/unidades.
% - Cajas tcolorbox temáticas para fundamentación técnica / teoría.
% - Bloques de código bash / programación con listings y resaltado.
% - Tablas compactas y estructuradas con tabularx / booktabs.
% =========================================================================

\documentclass[9pt,a4paper]{extarticle}
\usepackage[utf8]{inputenc}
\usepackage[T1]{fontenc}
\usepackage{lmodern}
\usepackage[spanish,es-tabla,es-noquoting]{babel}
\usepackage[top=1.2cm, bottom=1.3cm, left=1.25cm, right=1.25cm]{geometry}
\usepackage{microtype}
\usepackage{xcolor}
\usepackage{amsmath,amssymb}
\usepackage{booktabs}
\usepackage{tabularx}
\usepackage{enumitem}
\usepackage{fancyhdr}
\usepackage{titlesec}
\usepackage{listings}
\usepackage[many]{tcolorbox}
\usepackage{hyperref}

% --- Configuración de Página y Tipografía ---
\pagestyle{plain} % Sin encabezado superior, únicamente pie de página limpio
\setlength{\parindent}{0pt}
\setlength{\parskip}{3pt}

% --- Paleta de Colores Pastel Profesionales y Armónicos ---
\definecolor{bannerbg}{HTML}{EDF2F7}
\definecolor{bannerframe}{HTML}{4A5568}
\definecolor{bannertitle}{HTML}{1A202C}

% Unidad / Módulo 1 (Azul Acero Pastel)
\definecolor{tema1bg}{HTML}{EBF4FA}
\definecolor{tema1frame}{HTML}{2B6CB0}
\definecolor{tema1text}{HTML}{1A365D}

% Unidad / Módulo 2 (Verde Salvia Pastel)
\definecolor{tema2bg}{HTML}{EAF7ED}
\definecolor{tema2frame}{HTML}{2F855A}
\definecolor{tema2text}{HTML}{1C4532}

% Unidad / Módulo 3 (Melocotón / Ámbar Pastel)
\definecolor{tema3bg}{HTML}{FEF3E9}
\definecolor{tema3frame}{HTML}{C05621}
\definecolor{tema3text}{HTML}{652B19}

% Unidad / Módulo 4 (Lavanda / Púrpura Pastel)
\definecolor{tema4bg}{HTML}{F6EEFB}
\definecolor{tema4frame}{HTML}{6B46C1}
\definecolor{tema4text}{HTML}{322659}

% Cajas de Código y Bloques de Comandos
\definecolor{codebg}{HTML}{F8FAFC}
\definecolor{codeframe}{HTML}{CBD5E1}
\definecolor{codetitlebg}{HTML}{E2E8F0}
\definecolor{codetitle}{HTML}{1E293B}

\definecolor{cmdkeyword}{HTML}{0284C7}
\definecolor{cmdstring}{HTML}{B45309}
\definecolor{cmdcomment}{HTML}{64748B}

% --- Estilo de Listings para Código ---
\lstdefinestyle{bashstyle}{
    language=bash,
    basicstyle=\ttfamily\fontsize{8.5pt}{10.2pt}\selectfont\color{black},
    keywordstyle=\color{cmdkeyword}\bfseries,
    stringstyle=\color{cmdstring},
    commentstyle=\color{cmdcomment}\itshape,
    showstringspaces=false,
    breaklines=true,
    breakatwhitespace=false,
    tabsize=2,
    columns=fullflexible,
    keepspaces=true,
    aboveskip=0pt,
    belowskip=0pt
}

% --- Cajas TColorBox Personalizadas ---
\newtcolorbox{bannerbox}{
    enhanced,
    colback=bannerbg,
    colframe=bannerframe,
    arc=2.5mm,
    boxrule=0.9pt,
    left=9pt, right=9pt, top=6pt, bottom=6pt
}

\newtcolorbox{bloqueheader}[4]{
    enhanced,
    colback=#1,
    colframe=#2,
    coltitle=#3,
    arc=2mm,
    boxrule=1.1pt,
    title={\textbf{\fontsize{10.5pt}{12.5pt}\selectfont #4}},
    left=8pt, right=8pt, top=4pt, bottom=4pt,
    before skip=6pt, after skip=3pt
}

\newtcolorbox{conceptbox}[4][]{
    enhanced,
    breakable,
    colback=#2,
    colframe=#3,
    coltitle=#4,
    arc=1.5mm,
    boxrule=0.5pt,
    leftrule=3.5pt,
    title={\textbf{\fontsize{8.5pt}{10.5pt}\selectfont $\blacktriangleright$\ Fundamentación Teórica / Concepto Clave: #1}},
    left=7pt, right=7pt, top=3pt, bottom=3pt,
    before skip=2.5pt, after skip=3.5pt,
    fontupper=\fontsize{8.5pt}{10.6pt}\selectfont
}

\newtcolorbox{codebox}[2][]{
    enhanced,
    breakable,
    colback=codebg,
    colframe=codeframe,
    arc=1.5mm,
    boxrule=0.7pt,
    title={\textbf{\fontsize{8.5pt}{10.5pt}\selectfont #2}},
    coltitle=codetitle,
    colbacktitle=codetitlebg,
    left=6pt, right=6pt, top=3pt, bottom=3pt,
    before skip=2.5pt, after skip=2.5pt,
    #1
}

% --- Formato Limpio de Subsecciones (Sin Duplicación de Numeración) ---
\titleformat{\subsection}
  {\fontsize{9.2pt}{11.2pt}\bfseries\color{bannerframe!90!black}}
  {}{0pt}{}

\titlespacing*{\subsection}{0pt}{4pt}{2pt}

\hypersetup{
    colorlinks=true,
    linkcolor=blue!70!black,
    urlcolor=blue!70!black,
    citecolor=blue!70!black
}

\begin{document}

% =========================================================================
% ENCABEZADO COMPACTO DIRECTO (PÁGINA 1)
% =========================================================================
\begin{bannerbox}
    \begin{minipage}[c]{0.72\textwidth}
        {\textbf{\fontsize{12.5pt}{14.5pt}\selectfont\color{bannertitle} [TÍTULO DEL APUNTE O GUÍA TÉCNICA]}}\\[2pt]
        {\fontsize{8.5pt}{10.5pt}\selectfont\color{bannerframe} \textbf{[Subtítulo Descriptivo o Unidades Comprendidas]}}
    \end{minipage}
    \hfill
    \begin{minipage}[c]{0.25\textwidth}
        \raggedleft
        {\fontsize{7.5pt}{9.5pt}\selectfont\color{bannerframe}
        \textbf{Materia:} [Sigla / Nombre]\\
        \textbf{Institución:} FI - UNJu\\
        \textbf{Tipo:} Apunte Compacto}
    \end{minipage}
\end{bannerbox}

\vspace{1pt}

% =========================================================================
% MÓDULO / UNIDAD 1 (AZUL PASTEL)
% =========================================================================
\begin{bloqueheader}{tema1bg}{tema1frame}{tema1text}{Módulo 1: [Nombre de la Unidad / Tema]}
Resumen ejecutivo o descripción introductoria del módulo 1.
\end{bloqueheader}

\subsection{1.1 [Título del Subtema]}
Texto explicativo del subtema...

\begin{codebox}{[Título del Bloque de Código o Comandos]}
\begin{lstlisting}[style=bashstyle]
# Comentario en ASCII estándar
comando --opcion argumento
\end{lstlisting}
\end{codebox}

\begin{conceptbox}[[Título del Concepto o Fundamentación]]{tema1bg!60!white}{tema1frame}{tema1text}
\begin{itemize}[leftmargin=*,noitemsep,topsep=1pt]
    \item \textbf{Punto clave 1:} Explicación del concepto o por qué.
    \item \textbf{Punto clave 2:} Fundamento técnico detallado.
\end{itemize}
\end{conceptbox}

% =========================================================================
% MÓDULO / UNIDAD 2 (VERDE PASTEL)
% =========================================================================
\begin{bloqueheader}{tema2bg}{tema2frame}{tema2text}{Módulo 2: [Nombre de la Unidad / Tema]}
Resumen ejecutivo del módulo 2.
\end{bloqueheader}

\subsection{2.1 [Título del Subtema]}
Texto explicativo...

\begin{conceptbox}[[Título del Concepto]]{tema2bg!60!white}{tema2frame}{tema2text}
\begin{itemize}[leftmargin=*,noitemsep,topsep=1pt]
    \item \textbf{Detalle:} Explicación conceptual.
\end{itemize}
\end{conceptbox}

% =========================================================================
% RESUMEN / TABLA COMPARATIVA
% =========================================================================
\vspace{4pt}
\begin{bloqueheader}{bannerbg}{bannerframe}{bannertitle}{Resumen y Matriz Comparativa}
\begin{center}
\renewcommand{\arraystretch}{1.15}
\fontsize{8pt}{10pt}\selectfont
\begin{tabularx}{\textwidth}{l p{3cm} X}
\toprule
\textbf{Elemento} & \textbf{Identificador / Parámetro} & \textbf{Descripción / Función} \\
\midrule
\textbf{Item 1} & Valor A & Descripción detallada del item 1. \\
\midrule
\textbf{Item 2} & Valor B & Descripción detallada del item 2. \\
\bottomrule
\end{tabularx}
\end{center}
\end{bloqueheader}

\end{document}
